\section{Local Orbital Stability via the Shape Energy}
\label{sec:corotating}

Section~\ref{sec:second} showed that the second controller is handled by a Lyapunov function whose
derivative is a sum of negative squares. For the \emph{first} controller no such \emph{global}
monotone Lyapunov function exists: the repulsive potential is non-convex, and the rotational phase
is a neutral mode. What is available instead is a sharp \emph{local} result, built from the
second-order structure of the shape subsystem and a simple energy. Combined with the selection
theorem (Proposition~\ref{prop:omega}), this pins the steady state to $|\omega_s|=\sqrt{k_2}$.

\subsection{Second-order shape dynamics and the shape energy}
\label{sec:second_order}

Let $\txi_i^s=\txi_i-\bx$ and $\tvell_i^s=\tvell_i-\bv$ denote the zero-mean shape variables (the
average part obeys the exponentially stable oscillator \eqref{eq:osc} and decouples exactly).
Differentiating $\dot{\txi}_i^s=\phi_i-k_1\txi_i^s+\tvell_i^s$ and eliminating the observer through
$\dot{\tvell}_i^s=-k_2\txi_i^s$ yields the \emph{second-order} form
\begin{equation}
  \ddot{\txi}_i^s+k_1\dot{\txi}_i^s+k_2\txi_i^s=\dot{\phi}_i .
  \label{eq:second}
\end{equation}
Because $\phi$ is the gradient of the repulsive potential
$P=\sum_{\text{edges}}\Phi(\norm{x_{ij}})$ with $\Phi'=\alpha$ (so $\phi=\nabla_{\txi^s}P$), its time
derivative is $\dot{\phi}=\nabla^2 P\,\dot{\txi}^s$. Hence the shape subsystem is the closed
second-order system
\begin{equation}
  \boxed{\;\ddot{\txi}^s+B(\txi^s)\,\dot{\txi}^s+k_2\txi^s=0,\qquad
  B(\txi^s)=k_1I-\nabla^2 P(\txi^s),\;}
  \label{eq:secondB}
\end{equation}
in which all damping information is carried by the symmetric matrix $B$. Now consider the energy
\begin{equation}
  E=\frac12\norm{\dot{\txi}^s}^2+\frac{k_2}{2}\norm{\txi^s}^2 .
  \label{eq:energy}
\end{equation}
Along \eqref{eq:secondB},
\begin{equation}
  \dot E=\dot{\txi}^{s\mathsf T}\bigl(\ddot{\txi}^s+k_2\txi^s\bigr)
       =-\dot{\txi}^{s\mathsf T}B(\txi^s)\,\dot{\txi}^s .
  \label{eq:Edot}
\end{equation}
Thus $E$ is non-increasing precisely on the set where $B\succeq0$.

\subsection{Rotation covariance and the neutral rotational mode}
\label{sec:rotcov}

The repulsive force is rotation-covariant: rotating every relative position rotates every force,
$\phi(R\txi)=R\phi(\txi)$ for every rotation matrix $R$. Differentiating this identity along the
generator $J$ at the current configuration gives the exact relation
\begin{equation}
  \nabla^2 P(\txi^s)\,(J\txi^s)=J\phi(\txi^s),
  \label{eq:rotcov}
\end{equation}
and hence, by \eqref{eq:secondB},
\begin{equation}
  B\,(J\txi^s)=k_1J\txi^s-J\phi(\txi^s) .
  \label{eq:Brot}
\end{equation}
At a radial-balance configuration $\phi_i=k_1\txi_i^s$ the right-hand side vanishes identically:
\begin{equation}
  B\,(J\txi^s)=k_1J\txi^s-J(k_1\txi^s)=0 .
  \label{eq:Brot0}
\end{equation}
The rotational direction $J\txi^s$ is therefore \emph{always} a null mode of the damping matrix $B$
at equilibrium---an exact, symmetry-free statement, and the second-order counterpart of the neutral
rotational phase: a rigid rotation $\dot{\txi}^s=\omega J\txi^s$ gives $\dot E=0$ exactly.

\subsection{What the shape energy proves locally}
\label{sec:lasalle}

Suppose that at the equilibrium $B\succeq0$ on the shape space, with the rotational direction its
\emph{only} null mode. This holds for the paper's example: at the radial-balance equilibrium of
Section~\ref{sec:verifyrot} the shape-space eigenvalues of $\nabla^2 P$ are
$\{-15.17,-15.17,-11.14,-6.36,-3.05,-3.03,-3.03,-0.14,-0.14,\ 0.5\}$, so those of
$B=k_1I-\nabla^2P$ are $\{0,\ 0.64,\ 0.64,\ 3.53,\ 3.53,\ 3.55,\ 6.86,\ 11.64,\ 15.67,\ 15.67\}\succeq0$,
with the single zero mode aligned to machine precision with the rotational direction $J\txi^s$
(purely tangential). (Analytically, $B\succeq0$ is the condition that the radial-balance
configuration is a local maximum of $W=P-\tfrac{k_1}{2}\sum_i\norm{\txi_i^s}^2$, i.e.\ that every
``breathing'' mode is stable.)

This energy identity should be read with one caveat. Since $B$ has a zero eigenvalue at the
equilibrium, continuity alone does \emph{not} imply $B\succeq0$ throughout a whole neighborhood;
small perturbations may make the neutral eigenvalue slightly negative. Thus \eqref{eq:Edot} is not,
by itself, a complete nonlinear Lyapunov proof of local orbital stability. What it does prove is the
exact structure of the neutral direction: at radial balance the only direction in which the energy
does not dissipate is the rotational phase.

Indeed, if a trajectory stays in a region where $B(\txi^s)\succeq0$ and
$\ker B(\txi^s)=\operatorname{span}\{J\txi^s\}$, then LaSalle's invariance principle confines the
$\omega$-limit set to
\begin{equation}
  \dot E=0\quad\Longleftrightarrow\quad
  \dot{\txi}^s\in\ker B=\operatorname{span}\{J\txi^s\},
  \label{eq:Edot0}
\end{equation}
i.e.\ the motion is a rigid rotation $\dot{\txi}^s=\omega J\txi^s$ at constant radii. On such an
invariant set \eqref{eq:secondB} collapses to the undamped oscillator
\begin{equation}
  \ddot{\txi}^s+k_2\txi^s=0 .
  \label{eq:osc2}
\end{equation}
For $\dot{\txi}^s=\omega J\txi^s$ we have $\ddot{\txi}^s=\omega J\dot{\txi}^s=-\omega^2\txi^s$, so
\eqref{eq:osc2} forces $\omega^2=k_2$. This recovers Proposition~\ref{prop:omega} on the limit set
and explains why the only non-dissipative motion allowed by the shape energy is the rigid rotation
at $|\omega|=\sqrt{k_2}$. The actual local orbital stability claim is therefore established below
by the co-rotating-frame linearization, not by the energy identity alone.

\subsection{Co-rotating frame and shape dynamics}
\label{sec:shape}

Let $\omega=\sqrt{k_2}$, and let $\omega_s=\pm\omega$ denote the \emph{signed} angular velocity,
fixed by the initial condition (the closed loop is mirror-symmetric under $J\mapsto-J$, so we
analyze $\omega_s=+\omega$ without loss of generality). The zero-mean shape variables
\begin{equation}
  \txi_i^s=\txi_i-\bx,\qquad \tvell_i^s=\tvell_i-\bv
  \label{eq:shape}
\end{equation}
obey the closed, autonomous subsystem
\begin{equation}
  \dot{\txi}_i^s=\phi_i-k_1\txi_i^s+\tvell_i^s,\qquad
  \dot{\tvell}_i^s=-k_2\txi_i^s,
  \label{eq:shape_dyn}
\end{equation}
obtained by subtracting the average dynamics \eqref{eq:avgx}--\eqref{eq:avgv}. Because $\phi_i$
depends only on the differences $x_i-x_j$, it is translation-invariant, so the average
\eqref{eq:osc} and the shape \eqref{eq:shape_dyn} decouple exactly. Now pass to the frame co-rotating
at $\omega$,
\begin{equation}
  \xi_i=e^{-J\omega t}\txi_i^s,\qquad w_i=e^{-J\omega t}\tvell_i^s .
  \label{eq:corotating}
\end{equation}
Using the rotation-equivariance $\phi_i(e^{J\omega t}\xi)=e^{J\omega t}\phi_i(\xi)$, the shape
dynamics become the \emph{autonomous} system
\begin{equation}
  \dot{\xi}_i=\phi_i(\xi)-k_1\xi_i-\omega J\xi_i+w_i,\qquad
  \dot{w}_i=-k_2\xi_i-\omega Jw_i .
  \label{eq:shape_rot}
\end{equation}
A fixed point of \eqref{eq:shape_rot} is a constant pair $(\xi^*,w^*)$. The second equation of
\eqref{eq:shape_rot} gives, using $J^{-1}=-J$,
\begin{equation}
  w_i^*=\frac{k_2}{\omega}\,J\xi_i^* ,
  \label{eq:wstar}
\end{equation}
which holds for \emph{every} $k_2>0$ and so imposes no condition on $\omega$: the observer is simply
locked to the tangential velocity. Substituting \eqref{eq:wstar} into the first equation and
separating the radial and tangential components yields
\begin{equation}
  \phi_i(\xi^*)=k_1\xi_i^*,
  \qquad
  \Bigl(\frac{k_2}{\omega}-\omega\Bigr)J\xi_i^*=0 ,
  \label{eq:equilibrium}
\end{equation}
that is, radial force balance together with $\omega^2=k_2$ whenever $\xi^*\ne0$. It is therefore the
\emph{first} equation, not the observer equation, that pins the rate; this recovers
Proposition~\ref{prop:omega} within the co-rotating picture. In the lab frame \eqref{eq:equilibrium} is the rigid
rotation $\txi_i(t)=e^{J\omega t}\xi_i^*$ of Section~\ref{sec:rotation}; conversely,
Proposition~\ref{prop:omega} states that \emph{any} rigid rotation must have
$|\omega|=\sqrt{k_2}$.

\subsection{Linearization}
\label{sec:lyaprot}

Let $W(\xi)=\sum_{\text{edges}}\Phi(\norm{x_{ij}})-\tfrac{k_1}{2}\sum_i\norm{\xi_i}^2$ with
$\Phi'(s)=\alpha(s)$, so that $\nabla_i W=\phi_i-k_1\xi_i$ and the equilibrium is $\nabla W=0$. Write
$H=\nabla^2 W(\xi^*)$ (symmetric); note that $H=\nabla^2P-k_1I=-B$, so the condition $B\succeq0$ of Section~\ref{sec:lasalle} is exactly $H\preceq0$. Linearizing \eqref{eq:shape_rot} about $(\xi^*,w^*)$ in the state
$[\xi;w]$ gives
\begin{equation}
  DF=\begin{bmatrix} H-\omega J & I\\ -\omega^2 I & -\omega J\end{bmatrix},
  \qquad \omega^2=k_2 .
  \label{eq:DF}
\end{equation}
The neutral rotational phase is exact: the direction $(\delta,\epsilon)=(J\xi^*,-\omega\xi^*)$,
tangent to the family $\{(e^{J\theta}\xi^*,\ \omega Je^{J\theta}\xi^*)\}_{\theta}$ of rotations of
the equilibrium, satisfies $DF\,(J\xi^*,-\omega\xi^*)=0$. This uses the rotation invariance of $W$,
which gives $\nabla^2 W_{\mathrm{rep}}(\xi^*)\,J\xi^*=J\nabla W_{\mathrm{rep}}(\xi^*)
=J\phi(\xi^*)=k_1J\xi^*$, so that $HJ\xi^*=\nabla^2 W_{\mathrm{rep}}J\xi^*-k_1J\xi^*=0$.

\subsection{Orbital stability criterion}
\label{sec:local}

Writing the observer residual $\tilde\eta_i=w_i-\omega J\xi_i$ and eliminating $w$ turns
\eqref{eq:DF} into the equivalent pair
\begin{equation}
  \dot\delta=H\delta+\tilde\eta,\qquad
  \dot{\tilde\eta}=-\omega JH\,\delta-2\omega J\,\tilde\eta,
  \label{eq:lin2}
\end{equation}
whose characteristic polynomial is the matrix quadratic
\begin{equation}
  \det\bigl[\lambda^2 I+\lambda(2\omega J-H)-\omega JH\bigr]=0 .
  \label{eq:quad}
\end{equation}
We say the rotating formation is \emph{locally asymptotically (orbitally) stable} if every root of
\eqref{eq:quad}---equivalently every eigenvalue of \eqref{eq:DF}---lies in the open left half-plane,
except the single neutral phase $\lambda=0$.

It is worth recording why the obvious simplification is unavailable here. If $H$ and $J$ commuted,
\eqref{eq:quad} would decouple, on each joint eigenspace, into
\begin{equation}
  \lambda=\tfrac12\Bigl(h-2\mathrm{i}\omega\pm\sqrt{h^2-4\omega^2}\Bigr),
  \label{eq:commuting}
\end{equation}
for each eigenvalue $h$ of $H$, so that $h<0$ yields $\operatorname{Re}\lambda<0$ and $h>0$ is
unstable. But this commuting reduction cannot describe the situation of interest, for two independent
reasons. First, at $h=0$ formula \eqref{eq:commuting} returns \emph{two} roots, $\lambda=0$ and
$\lambda=-2\mathrm{i}\omega$, and the latter sits on the imaginary axis; a null direction of $H$ would
thus contribute a second marginal mode rather than a single neutral phase. Second, if $[H,J]=0$ then
$\ker H$ is $J$-invariant and therefore of even real dimension, so ``a single null direction'' is
impossible in the commuting case. The commuting case is consequently never asymptotically stable
modulo one phase, and it must not be used as a proxy for the stability criterion.

For a finite polygon $H$ and $J$ do \emph{not} commute---only the discrete rotational symmetry
survives---and both obstructions disappear: $\ker H$ is one-dimensional, spanned by $J\xi^*$, and no
spurious $-2\mathrm{i}\omega$ mode appears. Equation~\eqref{eq:quad} must therefore be checked
directly, which the next subsection does for the paper's parameters.

\subsection{Outcome of the spectral test}
\label{sec:verifyrot}

Criterion \eqref{eq:quad} was evaluated numerically at the paper's parameters; the computation, the
equilibrium data, and the $k_2$-sweep are reported in Section~\ref{S-sec:sim_spectrum} of the
companion. The outcome is the one the theory predicts: the linearization $DF$ has exactly one zero
eigenvalue, its eigenvector aligns with $(J\xi^*,-\omega\xi^*)$ to six digits, and the spectral
abscissa excluding that mode is strictly negative. The rigid rotation is therefore locally
asymptotically orbitally stable at those parameters, with the rotational phase as the only neutral
direction---exactly the structure Sections~\ref{sec:rotcov}--\ref{sec:lasalle} predict, now with the
sign of the remaining spectrum settled.

Two features of the numbers matter for the analysis. First, the identity $HJ\xi^*=0$ of
Section~\ref{sec:rotcov} holds to $1.3\times10^{-14}$, confirming the neutral mode is the exact
consequence of rotation covariance rather than a numerical coincidence. Second, the decay rate scales
with $\omega$: it is negative across a factor of $20$ in $k_2$, but shrinks as $k_2\to0$, so the
orbit's basin is entered ever more slowly for small gains. That scaling is what makes the slow
spin-up transients at small $k_2$ a property of the linearization rather than an artefact of
integration.

\subsection{The local--global boundary}
\label{sec:localglobal}

The energy $E$ of \eqref{eq:energy} is \emph{not} a global Lyapunov function, because $B$ is
sign-indefinite away from equilibrium. The tangential curvature of the repulsion,
$\Phi'(s)/s=\alpha(s)/s$, diverges as $s\to d$; whenever two vehicles come close enough, the matrix
$B=k_1I-\nabla^2P$ acquires a negative eigenvalue. Since $\dot E=-2\,\dot\txi^{s\mathsf T}B\dot\txi^s$
by \eqref{eq:Edot}, $E$ can then \emph{increase}, and along a transient that pulls the formation
inward from a wide initial ring it does: $\lambda_{\min}(B)$ starts at $k_1$ while the vehicles are
out of sensing range, goes negative during the inward collapse, and returns to zero---the rotational
null mode---only once radial balance is reached. The companion records this crossing quantitatively
for the baseline run (Section~\ref{S-sec:sim_observations}).

This is the concrete mechanism by which the rotation is a limit cycle rather than a damped
equilibrium: the tangential curvature of the repulsion acts as a transient \emph{anti-damping} that
pumps energy into the rotational mode, and the neutral rotational direction of $B$ is exactly where
that energy input stops. It also delimits what Section~\ref{sec:lasalle} can claim. The energy
argument certifies a neighbourhood of the rotating branch on which $B\succeq0$; it says nothing about
initial conditions outside that neighbourhood, and the sign change of $\lambda_{\min}(B)$ shows the
excluded region is not empty. Global convergence remains open, and Conjecture~\ref{conj:dist}
records the three observed long-term behaviors.
